\scanpage{59}
les \uline{op. diff.} de $M$ dans $N$ (rel à $S$),

Alors on trouve encore un foncteur \deleted{\unsure{les|des}} \unsure{de|en} $M$, mais \deleted{\unsure{ici|(si $i$}} non à valeurs ds les $P^i$-Modules, mais \inserted{(si $i\geq1$)} seulement ds les $P^{i-1}$-Modules, où on a considéré $P^i[M]$ comme $P^{i-1}$-Module par \uline{restriction des scalaires} à l'aide de $\alpha^{i-1}:P^{i-1}\longrightarrow P^i$. \inserted{On le note alors $Q^{i-1}[M]$ $\circledast$} En particulier, $\deleted{P}Q^0[M]$ est \deleted{\unsure{de|est}} un foncteur, en $M\in\operatorname{Ob}\operatorname{Mod}(.X/S)$ \inserted{[Modules sur $\underline O_X$, avec comme hom les op. diff. \deleted{d'ordre} rel. à $S$]}, à valeurs ds les pro-$\underline O_X$-Modules.
% margin note written sideways
\margin{et qu'on utilise comme morphismes \unsure{de|des} \deleted{\unsure{objets|Modules}} \unsure{$\mathbb N$-pro-objets|N-pro-objets} ceux de Artin--Riesz, \unsure{Supposant|exprimant} qu'on ne considère que des op. diff de degré borné, ce qui est autom. si $X$ qu. cpct].}

\deleted{[$P^0[M]=M$ \deleted{est} doit être considéré comme foncteur à valeurs ds $\mathrm{Ab}(X)$].}

Notons que l'hom $\alpha^{\cdot}:P^{\unsure{*|1}}\mathbin{-}(SP)^*=Q^*$ permet de considérer $Q^*$ comme \deleted{\unsure{module}} \inserted{algèbre} \deleted{\unsure{sur}} cosimplicial sur l'\unsure{anneau} \deleted{\unsure{semi|déplacement}}cosimplicial $P^*$, et que les opérations extérieures droites
\[
\boxed{\circledast\quad Q^{i-1}[M]=P^i\otimes_{\underline O_X} M}
\]
% archive sheet number: 58

